På schemalagd tid senast sista läsveckan i december ska du avlägga ett obligatoriskt muntligt prov för handledare eller kursansvarig. Du måste vara godkänt på alla laborationer för att få göra det muntliga provet. Syftet med provet är att kontrollera att du har godkänd förståelse för de begrepp som ingår i kursen. Du rekommenderas att förbereda dig noga inför provet, t.ex. genom att gå igenom grundläggande begrepp för varje kursmodul och repetera grundövningar och laborationer.

Provet sker som ett stickprov ur kursens innehåll. Handledare (eller kursansvarig) börjar med att ge dig några slumpvis valda frågor där du ombeds förklara några av de begrepp som ingår i kursen. Du får även uppdrag att skriva kod som liknar kursens övningar och förklara hur koden fungerar. Vid osäkerheter eller kunskapsluckor kan handledare/kursansvarig ställa utvalda följdfrågor för att identifiera områden som kräver att du tränar mer. 

Om det visar sig oklart huruvida du uppnått godkänd förståelse kan du behöva komplettera ditt muntliga prov. Kontakta kursansvarig för information om omprov.  

\subsection{Kontinuerlig muntaträning}

Med hjälp av Muntabot \url{https://fileadmin.cs.lth.se/pgk/muntabot/} ska du kontinuerligt, vecka för vecka, kontrollera att du har förstått veckans grundläggande begrepp och att du självständigt kan skapa kod som illustrerar veckans koncept, enbart med hjälp av papper, penna och snabbreferensen. 

I muntaträningen för varje vecka ingår att i dialog med handledare visa att du självständigt med papper och penna kan skriva kod, t.ex. med en funktion eller algoritm, som illustrerar veckans grundläggande koncept. Om det år något du tycker är speciellt svårt så får du hjälp på resurtid att göra en plan för vilka övningar du behöver träna mer på. 

Hur hänger de olika delmomenten ihop? Grundövningarna ger förståelse för de grundläggande delar som ingår i kursen medan laborationerna ger förståelse för hur delarna kan sättas samman till en helhet i ett större program. Extraövningarna finns till för de som behöver träna mer. Den kontinuerliga muntaträningen hjälper dig att verifiera att du har de kunskaper som krävs för att gå vidare. Varje vecka bygger vidare på de kunskaper du fick föregående vecka så det är därför extra vikigt att du verifierar att du kan använda dina kunskaper praktiskt inför fortsättningen. 
